\documentclass[10pt,a4paper]{amsart}
\usepackage[margin=19mm]{geometry}
\usepackage{amsmath,amssymb,mathtools}
\usepackage[scr=rsfs,cal=euler]{mathalfa}
\usepackage{xfrac}
\usepackage[unicode,hidelinks,bookmarks=false]{hyperref}
\newcommand{\ie}{i.e.\ }
\newcommand{\cf}{cf.\ }
\newcommand{\resp}{respectively\ }
\newcommand{\Subsection}{\subsection}
\providecommand{\nobreakdash}{\nobreak\hbox{-}}
\setlength{\emergencystretch}{2em}
\multlinegap0pt
\usepackage{mathtools}
\newtheorem{theorem}{Theorem}
\newtheorem{lemma}{Lemma}
\def \R {\mathbb R}
\renewcommand {\chi}{\operatorname{1}}
\title{New criteria on positive-definite distributions}
\author{J. Haddad}
\date{Source: arXiv:2511.09281v2 (CC BY 4.0)}
\begin{document}
\maketitle

\noindent\emph{Source and licence.} New criteria on positive-definite distributions. Version arXiv:2511.09281v2 (CC BY 4.0), distributed by arXiv under the Creative Commons Attribution 4.0 International licence, http://creativecommons.org/licenses/by/4.0/, as recorded in the arXiv licence metadata for this exact version and shown on the abstract page https://arxiv.org/abs/2511.09281v2. Full licence notice: \texttt{licenses/arxiv-2511.09281-CC-BY-4.0.txt}.\par\smallskip

\noindent\textit{Editorial scope.} Counted unit: the article's theorem thm\_decreasing\_dimn (source0.tex lines 128--137). The article's complete original proof of it is delivered with it (source0.tex lines 267--302, one \texttt{proof} environment of 36 lines, which the article places away from the statement). No mathematics is written, repaired or re-derived: the statement and the proof are the article's own lines, in the article's own notation, and no retained line is substituted or re-flowed.  Retained uncounted as the source-local context the proof needs: lines 187--202.  Nothing was omitted from any retained span, so the package carries no omission record.  No retained line contains a citation, so the wrapper carries no bibliography and no bibliography entry.  No retained line contains a figure environment, an image-inclusion command or drawing code, so no image is delivered and the package declares no asset dependency.  Editorial formatting only, in addition: an AMS \texttt{amsart} preamble that restates the article's own declarations and macros used by the retained lines, namely the environments \texttt{theorem}, \texttt{lemma} on separate counters, together with the standard packages those lines need.  The retained environments print in this document as Theorem 1, Lemma 1 in order of appearance, whereas the article prints the same items with the numbering of its own full document.  The complete native source of the article as distributed in the arXiv e-print for this exact version is bundled unchanged as \texttt{sources/188969/source0.tex}.
\begin{lemma}
	\label{lem_decreasing_dim1}
	Let $\varphi, \psi:\R \setminus \{0\} \to \R$ be two non-negative, even, measurable functions, such that $\varphi(x)$ and $|x|^{-1} \psi(x)$ are non-increasing for $x>0$.
	Assume any of the three following conditions are met:
	\begin{itemize}
		\item $\varphi, \psi \in L^p(\R)$ with $p \in [1,2]$.
		\item $\varphi$ is bounded with compact support and $\psi(x) \min\{1, |x|^{-1}\}$ is integrable.
		\item $ \varphi \in L^1(\R) \cap C^3(\R)$ and $\psi \in L^1_{loc}(\R)$.
	\end{itemize}
	Then 
	\begin{equation}
		\label{eq_positive_product}
		\int_{-\infty}^{\infty} \hat \varphi(x) \psi(x) dx = \int_{-\infty}^{\infty} \varphi(x) \hat \psi(x) dx \geq 0.
	\end{equation}

\end{lemma}

\begin{theorem}
	\label{thm_decreasing_dimn}
	Let $f:(0,\infty) \to \R$ be a non-negative non-increasing function.
	Assume any of the following two conditions are satisfied:
	\begin{itemize}
		\item $n \geq 3$ and $f(r) \min\{1, r\}$ is integrable.
		\item $n \geq 9$ and $r f(r)$ is locally integrable.
	\end{itemize}
	Then the function $|x|^{-n+2} f(|x|)$ is a positive-definite distribution.
\end{theorem}

\begin{proof}[Proof of Theorem \ref{thm_decreasing_dimn}]
	First let us consider the case where $f(r) \min\{1, r\}$ is an integrable function.
	Extend $f:\R \setminus\{0\} \to \R$ as an even function and consider an even non-negative test function $\delta$. By polar coordinates,
\begin{align}
	\int_{\R^n} |x|^{-n+2} f(|x|) \hat \delta(x) d x
	&= \int_{S^{n-1}} \int_0^\infty r f(r) \hat \delta(r v) d r d v \\
	&= \frac 12 \int_{-\infty}^\infty r f(r) \int_{S^{n-1}} \hat \delta(r v) d v d r \\
	%&= \frac 12 \int_{-\infty}^\infty \hat f(r) \int_{S^{n-1}}\mathcal R \delta (v,r) dv dr \\
	&= \frac 12 \int_{-\infty}^\infty r f(r) \mathcal F\left[ \int_{S^{n-1}}\mathcal R \delta (v,\cdot) dv \right](r) d r.
\end{align}
The integral of the Radon transform can be written as
	\begin{equation}
		\label{eq_integral_radon}
		\int_{S^{n-1}}\mathcal R \delta (v,r) dv = \kappa_{n-2} \int_{\R^n} \delta(x) ( 1 - (r/|x|)^2 )_+^{\frac{n-3}2} |x|^{-1} dx,
	\end{equation}
	where $\kappa_d$ is the $(d-1)$-volume of $S^{d-1}$.
	To see this, write the Radon transform as $\mathcal R \delta (v,r) = \frac{\partial}{\partial r} \int_{\R^n} \delta(x) \chi_{(x, v) \leq r} dx$, then
\begin{align}
	\int_{S^{n-1}}\mathcal R \delta (v,r) d v
	&= \int_{\R^n} \delta(x) \frac{\partial}{\partial r} \int_{S^{n-1}} \chi_{(x, v) \leq r} d v d x,
\end{align}
which by rotational invariance can be computed as
\begin{align}
	\int_{S^{n-1}} \chi_{(x, v) \leq r} d v
		&= \int_{S^{n-1}} \chi_{ (e_1, v) \leq r/|x|} d v \\
		&= \kappa_{n-2} \int_{-1}^{r/|x|} (1-s^2)^{\frac{n-3}2} d s
\end{align}
if $r \leq |x|$, and $\kappa_{n-1}$ otherwise.
Take derivative with respect to $r$ to obtain \eqref{eq_integral_radon}.

Finally we get
	\[\int_{\R^n} |x|^{-n+2} f(|x|) \hat \delta(x) dx = \kappa_{n-2} \int_{\R^n} \frac {\delta(x)}{|x|} \frac 12 \int_{-\infty}^\infty r f(r) \mathcal F[ ( 1 - (r/|x|)^2 )_+^{\frac{n-3}2}] dr dx.\]
	For fixed $x \in \R^n \setminus \{0\}$, the function $\varphi_x(r) = ( 1 - (r/|x|)^2 )_+^{\frac{n-3}2}$ is bounded with compact support, even and non-increasing for $n\geq 3$, and the theorem follows from Lemma \ref{lem_decreasing_dim1} with $\psi(r) = r f(r)$ under the second condition.

	For the second part, notice that for $n \geq 9$, we have $\varphi_x \in C^3(\R)$ and we can use the second part of Lemma \ref{lem_decreasing_dim1}.
\end{proof}

\end{document}
